\chapter{Extended AllSides Dataset}
\label{appendix:extended_dataset}

\begin{table}[ht]
\centering
\small 
\begin{tabular}{c c c c c c}
\toprule
\textbf{Split} & \textbf{Total} & \textbf{Left} & \textbf{Center} & \textbf{Right} \\
\midrule
 Train (orig) & 26590 & 8861 & 7488 & 10241 \\
 Train (ext.) & 36774 & 14325 & 9033 & 13416 \\
 Test (orig) & 1300  & 402 & 299 & 599 \\
 Test (ext.) & 11516 & 4998 & 2575 & 3943 \\
\bottomrule
\end{tabular}
\caption{Allsides Media Split Dataset}
\label{tab:media_split_dataset}
\end{table}

The AllSidesV2 dataset was created by augmenting the original AllSides Media split with additional articles sourced from the AllSides website. To maintain methodological consistency, we ensured all new articles from sources present in the original test set were included in the extended test set, while all other articles were added to the training set. This approach preserves the source-based separation between training and test data, preventing source leakage that could artificially inflate performance metrics. 

Data collection adhered to the following protocol:

\begin{enumerate}
    \item Article metadata (e.g. title, source, publication date, bias rating) was scraped from AllSides following their terms of service, which allow research use.
    \item Articles were then scraped only from source websites that allow research use in their terms of service.
    \item Each article was verified to contain full text content exceeding 100 tokens.
\end{enumerate}

\begin{table}[ht]
\centering
\small
\begin{tabular}{lrr}
\toprule
\textbf{Source} & \textbf{Updated} & \textbf{Original} \\
\midrule
CNN (Web News) & 2398 & 101 \\
BBC News & 1125 & 100 \\
Reuters & 950 & 100 \\
The Guardian & 788 & 100 \\
Associated Press & 737 & 98 \\
Salon & 705 & 98 \\
Reason & 617 & 93 \\
Newsmax & 583 & 88 \\
Breitbart News & 571 & 97 \\
CBN & 534 & 98 \\
The Daily Caller & 509 & 98 \\
Guest Writer (Right) & 471 & 20 \\
American Spectator & 298 & 83 \\
ABC News & 296 & 97 \\
Newsmax (News) & 258 & 1 \\
Guest Writer (Left) & 212 & 2 \\
ABC News (Online) & 164 & 2 \\
Guest Writer & 99 & 1 \\
Newsmax (News) & 66 & 9 \\
John Stossel & 44 & 2 \\
Newsmax (Opinion) & 35 & 2 \\
The American Spectator & 18 & 1 \\
AP Fact Check & 8 & 2 \\
Ann Coulter & 7 & 1 \\
Tucker Carlson & 7 & 1 \\
Brent Bozell & 6 & 1 \\
Ben Stein & 5 & 2 \\
Matt Welch & 3 & 1 \\
Dick Morris & 2 & 1 \\
\bottomrule
\end{tabular}
\caption{Sources in original vs.\ AllSidesV2 dataset}
\label{tab:sources_comparison}
\end{table}

The original dataset, as described in Section~\ref{sec:dataset}, contains a total of 27,890 articles in the training set and 1,300 articles in the test set. AllSidesV2 increases the training set to 36,774 articles and the test set to 9,830 articles. Table~\ref{tab:sources_comparison} presents the distribution of sources in the original versus extended dataset. Notable expansions include CNN (Web News) increasing from 101 to 2,398 articles, BBC News from 100 to 1,125, and Reuters from 100 to 950. The expanded dataset available on \href{https://huggingface.co/datasets/dragonslayer631/article-bias-prediction-media-splits-updated}{HuggingFace} under the Apache 2.0 licence. This is the same as the original dataset published by \citet{baly2020we}, which can be found on \href{https://github.com/ramybaly/Article-Bias-Prediction/tree/main?tab=readme-ov-file}{GitHub}.

\section{Classifier Performance by Evaluation Paradigm}
\label{appendix:leakage-stats}

The following tables report F1-macro and Cohen's $\kappa$ for the best-performing domain-adapted model ($T_{id}(32){+}T_h$) under the media-split and random-split evaluation paradigms described in Section~\ref{sec:source-leakage}. Media-split results are computed on the full AllSidesV2 test set ($n=9{,}830$); random-split results are computed on a 4,104-article subset in which sources overlap between training and test.

\begin{table}[ht]
\centering
\small
\begin{tabular}{lrrrrrr}
\toprule
 & \multicolumn{3}{c}{\textbf{Media Split}} & \multicolumn{3}{c}{\textbf{Random Split}} \\
\textbf{Year} & $n$ & F1-Macro & $\kappa$ & $n$ & F1-Macro & $\kappa$ \\
\midrule
2012 & 462  & 48.3 & 0.33 & 114 & 69.2 & 0.59 \\
2013 & 836  & 48.9 & 0.44 & 248 & 74.2 & 0.79 \\
2014 & 811  & 55.8 & 0.50 & 212 & 81.4 & 0.80 \\
2015 & 840  & 62.1 & 0.55 & 244 & 74.3 & 0.69 \\
2016 & 1060 & 60.0 & 0.58 & 325 & 71.4 & 0.67 \\
2017 & 1042 & 61.7 & 0.45 & 344 & 90.3 & 0.85 \\
2018 & 1296 & 56.2 & 0.35 & 387 & 90.4 & 0.86 \\
2019 & 1520 & 59.6 & 0.42 & 431 & 90.3 & 0.85 \\
2020 & 1018 & 55.9 & 0.37 & 501 & 85.4 & 0.78 \\
2021 & 189  & 54.1 & 0.38 & 281 & 85.9 & 0.81 \\
2022 & 253  & 57.4 & 0.33 & 379 & 87.0 & 0.83 \\
2023 & 242  & 55.7 & 0.29 & 333 & 86.9 & 0.81 \\
2024 & 206  & 55.1 & 0.32 & 253 & 87.2 & 0.81 \\
2025 & 46   & 57.0 & 0.56 & 48  & 92.0 & 0.87 \\
\bottomrule
\end{tabular}
\caption{Year-by-year F1-macro and Cohen's $\kappa$ under each evaluation paradigm. The leakage gap widens sharply from 2017 onward as the random split improves to over 90\% while the media split remains in the 54--62\% range, confirming a source-routing explanation rather than a content-difficulty effect. The 2025 rows have small $n$ and should be interpreted with caution.}
\label{tab:year_leakage}
\end{table}

\begin{table}[ht]
\centering
\small
\begin{tabular}{lcrrrr}
\toprule
 & & \multicolumn{2}{c}{\textbf{Media Split}} & \multicolumn{2}{c}{\textbf{Random Split}} \\
\textbf{Source} & \textbf{Label} & $n$ & F1-Macro & $n$ & F1-Macro \\
\midrule
ABC News           & Left   & 197  & 10.3 & 40  & 26.9 \\
Guest Writer       & Mixed  & 98   & 14.8 & 24  & 18.2 \\
Reason             & Right  & 524  & 16.9 & 116 & 58.5 \\
Reuters            & Center & 842  & 17.4 & 187 & 33.0 \\
Newsmax (News)     & Right  & 207  & 18.1 & 47  & 29.4 \\
BBC News           & Center & 1023 & 18.5 & 221 & 32.3 \\
Newsmax            & Right  & 489  & 18.5 & 114 & 28.3 \\
Newsmax (News alt.)& Right  & 56   & 18.8 & --- & ---  \\
CBN                & Right  & 432  & 23.1 & 105 & 32.2 \\
ABC News (Online)  & Left   & 147  & 26.3 & 33  & 32.8 \\
Salon              & Left   & 606  & 26.4 & 137 & 30.1 \\
The Daily Caller   & Right  & 206  & 29.3 & 57  & 32.7 \\
The Guardian       & Left   & 686  & 30.9 & 143 & 32.4 \\
Guest Writer (Left)& Left   & 205  & 31.3 & 47  & 31.5 \\
Breitbart News     & Right  & 434  & 31.4 & 92  & 32.0 \\
Guest Writer (Right)& Right & 445  & 32.1 & 114 & 32.4 \\
American Spectator & Right  & 215  & 32.5 & 46  & 32.6 \\
CNN (Web News)     & Center & 2283 & 33.0 & 518 & 33.1 \\
Associated Press   & Center & 629  & 45.2 & 147 & 63.6 \\
\bottomrule
\end{tabular}
\caption{Per-source F1-macro under each evaluation paradigm for sources with at least 20 articles in the media-split test set, sorted by media-split F1-macro. ``---'' indicates the source has fewer than 20 articles in the random-split subset. For single-class sources, macro F1 has a ceiling of approximately 33\%; performance below this ceiling reflects misclassification of the source's articles. Associated Press appears last as it has articles in multiple ideology classes. The Label column reflects AllSides political ratings.}
\label{tab:source_leakage}
\end{table}

\begin{table}[ht]
\centering
\small
\begin{tabular}{lrrrr}
\toprule
 & \multicolumn{2}{c}{\textbf{Media Split}} & \multicolumn{2}{c}{\textbf{Random Split}} \\
\textbf{Topic} & $n$ & Miss (\%) & $n$ & Miss (\%) \\
\midrule
\multicolumn{5}{l}{\textit{Source-leakage-dependent topics}} \\
disaster            & 29 & 89.7 & --- & --- \\
ukraine\_war        & 12 & 66.7 & --- & --- \\
us\_military        & 46 & 56.5 & 9   & 22.2 \\
labor               & 18 & 55.6 & 5   & 60.0 \\
sports              & 40 & 55.0 & --- & --- \\
business            & 34 & 52.9 & --- & --- \\
sexual\_misconduct  & 45 & 51.1 & --- & --- \\
\midrule
\multicolumn{5}{l}{\textit{Content-level hard topics (persistent regardless of leakage)}} \\
race\_and\_racism   & --- & $<$50 & 34 & 38.2 \\
polarization        & --- & $<$50 & 66 & 34.8 \\
economic\_policy    & --- & $<$50 & 15 & 26.7 \\
terrorism           & --- & $<$50 & 37 & 24.3 \\
\bottomrule
\end{tabular}
\caption{Topic-level misclassification rates under each evaluation paradigm (topics with $n \geq 5$ in the respective split). ``---'' indicates the topic did not appear in that split's test set or had fewer than 5 articles. ``$<$50'' indicates the topic was not among the 20 worst-performing topics in that split, implying a miss rate below approximately 50\%. Source-leakage-dependent topics show high miss rates in the media split but low miss rates (or absence) in the random split; content-level hard topics remain difficult in the random split despite resolved source routing.}
\label{tab:topic_leakage}
\end{table}
